Ao observar, registrar ou extrair uma amostra de dados a partir de um fenômeno aleatório,
isto é, cujo resultado não pode ser deterministicamente previsto, a Estatística, dentre
outras coisas, procura entender o comportamento das frequências de ocorrência dos valores
amostrados.
Essas frequências são estimativas das \textbf{probabilidades} dos eventos observados.
A partir de hipóteses adequadas, podemos modelar as frequências por modelos probabilísticos
que, por sua vez, são objetos de estudo da Teoria das Probabilidades.

Probabilidades são frequentemente representadas como porcentagens e são muitas vezes
entendidas como a frequência de um evento quando se repete o experimento um número
suficiente de vezes. Por exemplo, dizer que 16\% da população possui diploma de
graduação significaria que se escolhermos aleatoriamente 100 pessoas dessa população,
cerca de 16 delas teriam diploma teriam diploma de graduação. Essa abordagem é chamada
de \textbf{frequentista}, sendo a mais amplamente adotada em aplicações estatísticas
clássicas.

Outra abordagem para compreender probabilidades é enquanto o grau de crença numa hipótese,
como ao dizer que se tem 80\% de certeza em dizer que irá chover, visto as evidências do
céu nublado e ventos frios. Essa é uma visão mais intuitiva e moderna de probabilidades,
sendo chamada de \textbf{Bayesiana}, sendo a base da Inferência Bayesiana.

A moderna Teoria das Probabilidades é fundamenta utilizando-se conceitos da Teoria
da Medida, um ramo da Análise que generaliza as noções de comprimento, área e volume
para espaços abstratos associando-os a uma ``medida'', a partir da qual também
estende-se uma Teoria de Integração para esses espaços. Probabilidades são vistas,
portanto, como medidas das chances de ocorrência de um subconjunto de eventos
dentro de um universo de resultados possíveis.
